\documentclass[a4paper]{article}
\usepackage[pdftex]{graphicx}
\usepackage[utf8]{inputenc}
\usepackage{enumerate}
\usepackage{amssymb}
\usepackage{icomma}
\usepackage{siunitx}
\sisetup{locale=DE}
\usepackage{tikz}
\usepackage{href-ul}
\hypersetup{
	colorlinks=true,
	linkcolor=blue,
	urlcolor=blue}
\usepackage{geometry}
\geometry{a4paper, top=15mm, left=15mm, right=15mm, bottom=15mm,
	headsep=10mm, footskip=12mm}

\begin{document}
	\begin{tikzpicture}(10,3)
		\draw[ultra thick](2,0) --(10,0) -- (10,1.3) --(2,1.3) -- (2,0);
		\draw[fill=black](2,0)-- (0,.6) -- (2,.6) -- (2,0);
		\node at (6,.6) {\large solution sheet from \href{http://www.okuyakl.com}{www.okuyakl.com}};
	\end{tikzpicture}
	\vspace{0.5 cm}
	
	\noindent{\bf Task 1. a)}\\
	
	\begin{minipage}{0.6\textwidth}
		Construction of the angle bisector:
		\begin{itemize}
			\item Circle k(A;r) around A with arbitrary radius
			\item Mark the intersection points of k on the legs [AD] and [AC].
			\item Intersect circular arcs around the points
			\item Straight line through A and the intersection of the circular arcs = angle bisector
		\end{itemize}
		
		\noindent{\bf Task 1. b)}\\
		Construction of the rectangle:
		\begin{itemize}
			\item distance $\overline{DC}$ with the compass
			\item circular arc with this radius around A
		    \item distance $\overline{DA}$ with the compass
			\item circular arc with this radius around C
			\item intersect both arcs = B
			\item Draw route [AB] and [CB].
		\end{itemize}
		
		\noindent{\bf Task 1. c)}\\
		Construction of the mirror point D':
		\begin{itemize}
			\item Intersect circular arc around D with axis AC twice
			\item Circular arcs around the two intersection points
			\item Straight line through intersections of circular arcs = perpendicular from D to [AC]
			\item Length of point D -- Remove the plumb base point with the compass
			\item circular arc with this radius intersects around L with straight lines on the other side = D'
		\end{itemize}
	\end{minipage}
	\begin{minipage}{0.4\textwidth}
		\includegraphics[width= 6 cm]{dreivier249}
		\includegraphics[width= 4 cm]{dreisp249}
	\end{minipage}
	
	\noindent{\bf Task 2.}\\
	\begin{minipage}{0.4\textwidth}
		\begin{itemize}
			\item Circles around A and B with any, equal radius
			\item Connecting intersections of circles results in perpendicular bisector $m_{[AB]}$.
			\item Measure the distance between 16 m minus 5 m = 11 m with a compass.
			\item Mark 11 m on the vertical bisector results in a penalty point.
		\end{itemize}
	\end{minipage}
	\begin{minipage}{0.6\textwidth}
		\includegraphics[width=8 cm]{torelf249}\\
	\end{minipage}
	
	\noindent{\bf Task 3.}\\
	A triangle is uniquely determined by two sides and the angle opposite the larger side.
	
	\noindent{\bf Task 4.}\\
	
	\noindent
	\begin{minipage}{0.7\textwidth}
		The triangle $\Delta A_2 B_2 C_2 $ emerges from the triangle $ \Delta A_1 B_1 C_1 $ by rotating it clockwise. Therefore both triangles are congruent to each other.
	\end{minipage}
	\hfill
	\begin{minipage}{0.3\textwidth}
		\includegraphics[width=5 cm]{pladrei053}
	\end{minipage}
	
	\noindent{\bf Task 5. a)}\\
	This triangle cannot be constructed because $\gamma$, the largest angle, is not opposite the largest side.
	
	\noindent{\bf Task 5. b)}\\
	It can be constructed because the triangle inequality applies
	$$ a+b > c; $$
	Construction description:
	\begin{itemize}
		\item Drawing the route [$AB$] = c
		\item circle around $A$ with radius $\SI{6,5}{\centi\meter}$
		\item circle around $B$ with radius $\SI{3,7}{\centi\meter}$
		\item intersection of both circles = $C$
		\item Drawing the routes [$AC$] and [$BC$]
	\end{itemize}
	
	\noindent{\bf Task 6.}\\
	The bisector is constructed for each triangle angle. These intersect at a point $S$. From this point you drop the perpendicular $L$ onto one side of the triangle. We draw a circle around $S$ with radius $\overline{SP}$. This is the incircle.
	\vspace{0.5 cm}
	
	
	\noindent
	\begin{minipage}{0.5\textwidth}
		\noindent{\bf Task 7. a)}\\
		\includegraphics[width=6 cm]{drav053}
	\end{minipage}
	\hfill
	\begin{minipage}{0.5\textwidth}
		\noindent{\bf Task 7. b)}\\
		\includegraphics[width=6 cm]{gtra053}
	\end{minipage}
	
	\noindent{\bf Task 8.}\\
	
	\begin{enumerate}[a)]
		\item Correct
		\item False. Counterexample: $\alpha=\beta=\gamma=80^\circ \qquad \delta=120^\circ$
		\item False. Counterexample: A parallelogram is not axisymmetric.
		\item False. Counterexample: An isosceles trapezoid is not point-symmetrical.
	\end{enumerate}
	
	\newpage
	\noindent{\bf Task 9.}\\
	
	\includegraphics[width=14 cm]{neun053}
	
	c) [$DA$] passes through the second and third quadrants.\\
	
	d) $d(B;g_{MC})=\SI{1,6}{\centi\meter}$\\
	
	f) $D'(4|-1)$; lies on the circle around M and on the extension of the line [$BC$].
	
	\begin{center}
		\includegraphics[width=7 cm]{../../viecher/eendcomic.pdf}
		
		Here you can go back to the \href{https://www.okuyakl.de/math/m8dreiL053/ae053.pdf}{task sheet}
	\end{center}
	
\end{document}